\documentclass{report}
\usepackage{amsmath,amssymb}
\setlength{\parindent}{0mm}
\setlength{\parskip}{1em}
\begin{document}
\begin{center}
\rule{6in}{1pt} \
{\large Nikolaos, N M Missirlis \\
{\bf Preconditioning augmented linear systems}}

Department of Informatics and Telecommunications \\ University of Athens \\ Panepistimiopolis \\ 15784 \\ Ilisia \\ Athens \\ Greece
\\
{\tt nmis@di.uoa.gr}\\
Maria Louka\end{center}

Let $A \in \mathbb{R}^{m\times m}$ be a symmetric positive definite
matrix and $B\in\mathbb{R}^{m\times n}$ be a matrix of full column rank,
where $m\geq n$. Then, the augmented linear system is of the form

\begin{equation}
\tilde Au=b\label{eq:01}
\end{equation}
where
\begin{equation}\label{eq:02}
\tilde A=\left(
\begin{array}{cc}
A & B\\
-B^{T} & 0
\end{array}
\right), \;\; u=
\left(
\begin{array}{cc}
x\\
y
\end{array}
\right),\;\; b=
\left(
\begin{array}{cc}
b_1\\
-b_2
\end{array}
\right)
\end{equation}
with $B^T$ be the transpose of the matrix B.
When A and B are large and sparse, iterative methods for solving
(\ref{eq:01}) are effective and more attractive than direct methods,
because of storage requirements and preservation of sparsity.
\\ \noindent The difficulty in applying iterative methods such as the
successive overrelaxation (SOR) method \cite{Young} to system
(\ref{eq:01}) is the singularity of the block diagonal part of the
coefficient matrix. Some methods have been developed to overcome this
difficulty such as the Uzawa and the Preconditioned Uzawa method
\cite{ArHurUz58}. In 2001, Golub et al. \cite{GolWuYuan2001} generalized
the Uzawa and the Preconditioned Uzawa method for the augmented system
(\ref{eq:01}) by introducing an additional acceleration parameter and
produced the SOR-like method. Other known methods, competitive and
specially attractive for solving (\ref{eq:01}) are the preconditioned
MINRES and GMRES methods \cite{FiRaSiWa98}. In case where the matrix A is
symmetric and positive definite, the Preconditioned Conjugate Gradient
(PCG) method \cite{HestSt52} is the most suitable. In 2003, Li, Evans and
Zhang \cite{LiEvansZha2003} proved that the PCG method is at least as
accurate as the SOR-like method. In 2005, Bai et al.
\cite{BaiParlWang2005} proposed the Generalized SOR (GSOR) method and
proved that it has the same rate of convergence but less complexity than
the PCG method. Furthermore, the Generalized Modified Extrapolated SOR
(GMESOR) method, a generalization of the GSOR method, was also proposed
for further study.
The latter is a generalization of the GSOR method as it uses one additional parameter. \\
\noindent In this paper, we study the Generalized Modified Extrapolated
SOR (GMESOR) method and the Generalized Modified Preconditioned
Simultaneous Displacement (GMPSD) method. The latter is derived by using
as preconditioning matrix the product of a lower and upper triangular
part of $\tilde A$ \cite{Louka}, \cite{LoukMis2}.
In an attempt to generalize our theory, we introduce an additional
parameter $a$ in the structure of the matrix $\tilde A$, with the hope of
increasing the rate of convergence for the aforementioned methods. Our
starting point, for studying these methods, is the derivation of
functional relations which relate the eigenvalues of their iteration
matrices with those of the key matrix $J=Q^{-1}B^TA^{-1}B$. Assuming that
the matrix Q is symmetric positive or negative definite matrix, the
eigenvalues of the matrix $J$ are real and either positive if the matrix
Q is positive definite or negative if the matrix Q is negative definite.
Under these assumptions we find sufficient conditions for the convergence
of the aforementioned methods and determine the optimum values of their
parameters. For the determination of the optimum values of the parameters
in GMESOR and GMPSD methods we followed a similar approach to Varga
\cite{Var1} p. 110, for the optimum $\omega$ of SOR.



\begin{thebibliography}{4}

\bibitem{ArHurUz58} K. Arrow, L. Hurwicz and H. Uzawa, Studies in
Nonlinear Programming, Stanford University Press, Stanford, 1958.

\bibitem{BaiParlWang2005} Z.-Z. Bai, B. N. Parlett and Z.-Q. Wang, On
generalized succesive overrelaxation methods for augmented linear
systems, Numer. Math. 102, (2005), 1-38.




\bibitem{FiRaSiWa98} B. Fischer, R. Ramage, D. J. Silvester, A. J.
Wathen, Minimum residual methods for augmented systems, BIT 38, (1998),
527-543.


\bibitem{GolWuYuan2001} G. H. Golub, X. Wu and J.-Y Yuan, SOR-like
methods for augmented systems, BIT 41, (2001), 71-85.


\bibitem{HestSt52} M. R. Hestenes and E. Stiefel, Methods of Conjugate
Gradients for Solving Linear Systems, Journal of Research of the National
Bureau of Standards, vol. 49, No. 6, (1952), 409-436.


\bibitem{LiEvansZha2003} C.-J. Li, Z. Li, D. J. Evans and T. Zhang, A
note on an SOR-like method for augmented
systems, IMA J. Numer. Anal. 23, (2003), 581-592.

\bibitem{Louka} M. A. Louka, Iterative Methods for solving linear
systems, Ph.D thesis (in Greek), 2011.

\bibitem{LoukMis2} M. A. Louka, N. M. Missirlis and F. I. Tzaferis, Is
modified PSD equivalent to modified SOR for two-cyclic matrices? Lin.
Alg. and its Appl., Vol. 432, No. 11, (2010), 2798-2815.


\bibitem{Var1} R. S. Varga, Matrix Iterative Analysis, Prentice-Hall, Inc.
Englewood Cliffs, N.J., 1962.





\bibitem{Young} D. M. Young, Iterative Solution of Large Linear Systems,
Academic Press, New York, 1971.

\end{thebibliography}


\end{document}
